% ------------------------------------------------------------------------------
% §5 Empirical design — drafted Phase C from:
%   Direction_clean/README.md (identification logic); 01_outcome findings;
%   stage2_gate_report.md; 03 findings.md (spec, deviations, inference);
%   04 findings.md (matching, power); findings_job2_contamination.md (clean-day rule)
% ------------------------------------------------------------------------------
\section{Empirical Design}\label{sec:design}

\subsection{The outcome: cheap capacity as a share of registered capacity}

The conduct of interest is how much capacity a unit makes available to the
market at prices at which it could actually run. The outcome is the
\emph{cheap-capacity share}: the megawatts a unit offers at or below a cheap
threshold, capped at its declared availability, divided by its registered
capacity. Two co-primary thresholds guard against threshold choice doing the
work: a fixed \$300/MWh, and a cost-indexed $2\times$ the unit's engineering
SRMC that month. The two definitions classify 96--100\% of intervals
identically per unit (worst case 3.9\%, Pelican Point), and a sweep from \$150
to \$500 and from $1.5\times$ to $3\times$ SRMC leaves the Torrens units'
withheld-interval share in an 82--96\% band throughout. Higher values of the
share mean \emph{less} withholding; a negative coefficient on any driver means
more withholding.

The share can fall through two channels with different regulatory treatment:
\emph{capacity withdrawn} (declared availability cut below normal --- physical
withholding) and \emph{capacity priced out} (availability normal, price above
the threshold --- economic withholding). Physical withdrawal dominates for
every unit (51--99\% of withheld intervals, rising to 91\% for Pelican Point),
with a material secondary share (10--33\%) doing both. The Torrens units'
unconditional distribution is bimodal: they sit at the zero-cheap floor in
69--77\% of \emph{all} intervals. Withholding is these units' resting state;
the empirical questions are whether the state deepens when the unit is
essential (RQ1) and whether it tracks the payment for being essential (RQ2).
Appendix~\ref{app:markup} reports the literature's markup lens on the same
offers as a cross-check on this outcome choice.

Three mechanical facts about this market discipline what the share can and
cannot mean, and Section~\ref{sec:results} exploits them directly. First,
the stance is indivisible: dispatch clears the cheapest offered megawatts
first, and a steam unit cannot produce its hundredth megawatt without its
first forty, so there is no bid in which the upper capacity is reachable
while the minimum-stable block is not --- any cheap tranche anywhere on the
ladder self-commits the whole unit. Second, self-commitment is what a
direction pays \emph{against}: compensation covers directed output net of
the counterfactual the unit's own bid establishes, so a unit whose bid
offers nothing has a zero counterfactual and is paid gross, while the same
megawatts offered cheaply would have run at (negative) spot. The one
payoff-relevant object in the bid is therefore binary --- is the
minimum-stable quantum within dispatch's reach, or not --- and the offered
\emph{price} is a dead instrument away from that threshold, which is
exactly how the units treat it (band prices change twice in roughly a
thousand observed bid revisions). Third, the share is accordingly best read
as a posture ruler: its zero point is the order-only stance (the unit runs
only if directed), its distance from zero measures how much running the
unit volunteers, and Section~\ref{sec:results} decomposes it into the
eligibility threshold and the depth beyond it.

\subsection{Essentiality: a rivals-only flag}

A unit-interval is \emph{essential} ($\pex=1$) when no combination of the
\emph{other} available South Australian synchronous units satisfies the minimum
synchronous requirement --- the system needs this unit (for Torrens, the
station) to complete a compliant combination. The flag is built exclusively
from rivals' declared availability; the unit's own bidding and availability
never enter. Two audits confirm the construction. A leakage audit regressing
the flag on the unit's own availability and cheap capacity yields
$R^{2}\in[0.0000, 0.0028]$: the flag is not a repackaging of the unit's own
behaviour. A degeneracy audit confirms essential intervals are non-trivial:
2--6\% of Torrens's and 66\% of Pelican Point's essential intervals still show
the unit bidding as usual, so essential does not mean withheld by
construction.

Essentiality is classified on realised system state rather than the
generator's bid-time forecast. Misclassification relative to the desk's
information set attenuates every essentiality coefficient toward zero; the
estimates below are lower bounds in this specific sense.\footnote{The
direction of the bias follows from the structure of the error, not from
assertion: the flag proxies the desk's expectation, and errors in it are
nondifferential with respect to the outcome conditional on the desk's
information --- realised-state surprises (a rival tripping after bids
closed) are uncorrelated with a bid formed before them. Classical
nondifferential misclassification of a binary regressor biases its
coefficient, and interactions with it, toward zero. What would overturn
this is misclassification correlated with the outcome through a third
channel; the leakage audit ($R^{2}\leq 0.0028$ on the unit's own behaviour)
bounds the obvious candidate.}

\subsection{The competition control: residual-demand slope}

Essentiality bundles two things: eligibility for direction compensation, and
whatever ordinary market power comes with being needed. To separate them, every
specification carries a direct measure of energy-market competition: the slope
of the unit's residual demand curve near the clearing price, built from rivals'
offer stacks net of interconnector imports. More negative slope means rivals
supply more at nearby prices --- more competition; a slope of exactly zero
means rivals are locally saturated --- maximum local market power. The zero
point is a mass point (6.7\% of Torrens intervals), so it enters as a separate
\emph{saturated} indicator alongside the continuous slope.

Two facts about this control, established before any regression, do
substantial interpretive work. First, competition and essentiality are nearly
orthogonal (correlation $-0.036$ for Torrens, $-0.018$ pooled): the control
cannot mechanically absorb the essentiality response. Second, their conditional
relationship runs \emph{opposite} to the market-power story: Torrens's
essential intervals average a residual-demand slope of $-5.04$~MW per \$/MWh
against $-3.31$ in non-essential intervals. Essential periods carry more rival
supply near the clearing price, not less, because essentiality is a
system-security condition that fires in high-renewables, low-demand periods ---
cheap-energy periods --- whereas energy-market scarcity fires elsewhere.
Section~\ref{sec:results} returns to this as the wrong-sign-of-conditions
result.

\subsection{Specifications}

\emph{RQ1.} The essentiality regression, at the unit-interval grain on the
four test units (1,261,576 rows),
\begin{equation}\label{eq:rq1}
\cheap_{it} \;=\; \beta\,\pex_{it} \;+\; \gamma_{1}\,\mathit{slope}_{it}
\;+\; \gamma_{2}\,\mathit{sat}_{it} \;+\; X_{it}'\delta \;+\; \alpha_{i}
\;+\; \mu_{m(t)} \;+\; \varepsilon_{it},
\end{equation}
with unit and month fixed effects and controls $X$ for SRMC, regional demand,
non-synchronous generation, and the spot price.\footnote{The contemporaneous
spot price is itself downstream of conduct and is a bad control in the
strict sense; it is retained because the committed specification included
it. Its own coefficient is indistinguishable from zero in every estimated
model ($p=0.77$--$0.96$ across Table~\ref{tab:rq1}'s specifications), so it
is not doing quiet work.}\footnote{The specification
originally called for the non-synchronous \emph{share} of demand; South
Australian regional demand crosses zero in-sample (rooftop solar), the share
explodes at the crossing, and with the broken control the essentiality
coefficients were inflated roughly tenfold. The level enters instead, with the
share version as a robustness row restricted to demand above 500~MW. The
deviation and its consequences are documented in the analysis record.}
Specification M1 omits the competition terms, M2 adds the slope, M3 adds the
saturated indicator; the movement of $\hat\beta$ from M1 to M3 measures how
much of the essentiality response the energy-market competition channel
absorbs.

\emph{RQ2.} The identifying test. Ordinary market power does not care what
the compensation formula pays; payment-seeking does. On a
coarsened-exact-matched (CEM) sample --- essential unit-intervals matched to
comparison intervals within unit $\times$ month $\times$ non-synchronous
quintile $\times$ hour block $\times$ competition bin cells (140,259 rows;
12,513 essential, a 100\% match rate) --- the regression adds the interaction
of essentiality with the compensation price,
\begin{equation}\label{eq:rq2}
\cheap_{it} \;=\; \theta\,\bigl(\pex_{it}\times \dt/100\bigr) \;+\;
\beta\,\pex_{it} \;+\; \gamma' C_{it} \;+\; X_{it}'\delta \;+\; \alpha_{i}
\;+\; \mu_{m(t)} \;+\; \varepsilon_{it}.
\end{equation}
$\dt$ is monthly, so month effects absorb its level; $\theta$ is identified
from cross-month variation in the essential-versus-matched gap. A negative
$\theta$ --- the gap in the cheap share widening as the payment rises --- is
the payment-seeking signature; a null is the insurance signature, conduct
triggered by the essential state but indifferent to what it pays. This mapping,
including the reading of a null as an informative bound, was written and
committed before estimation.

\emph{June 2022.} The east-coast crisis month contains a market suspension
window and an administered-price period, and holds 19\% of Torrens's essential
unit-rows. Every RQ2 estimate is reported under four treatments: excluding the
suspension window (base), excluding all of June, including the window at the
administered \$300 cap, and excluding pre-suspension June.

\emph{The estimand, stated precisely.} Because $\dt$ is monthly and the
mechanism evidence of Section~\ref{sec:mechanism} shows the posture does not
flex daily, $\theta$ is not a daily behavioural response. It is the
cross-month gradient of a standing posture's \emph{coverage}: how much more
of the essential intervals the order-only stance blankets, relative to
matched intervals, in months when a direction pays more. A policy chosen at
the operating-calendar frequency and tilted --- however it is decided ---
toward the payment environment produces exactly this object, and no claim
about interval-by-interval calibration is made or needed.

\subsection{Inference and hygiene}

Standard errors cluster by month. With 36 clusters (and effectively 21
essential-bearing months in RQ2), analytic clustered inference is supplemented
throughout by a wild cluster bootstrap ($R=999$; Rademacher weights primary,
Webb weights as sensitivity).\footnote{The bootstrap is the unrestricted
variant (\texttt{sandwich::vcovBS}); the null-imposed implementation could not
be installed in the analysis environment. With few effective and unbalanced
clusters the unrestricted variant can be liberal, so the headline estimate
is additionally subjected to randomization inference that imposes the null
by construction (Section~\ref{sec:robustness}), and the null-imposed
bootstrap proper remains a pre-submission upgrade. Analytic and bootstrap
$p$-values agree closely throughout, and every fixed-effects estimate was
re-estimated by explicit-dummy least squares and confirmed numerically
identical.} RQ2's power structure is reported with the
result, not discovered after it: the essential mass spans 21 month-clusters
with 50.8\% of it in the top three, and the compensation price varies across
essential rows with a standard deviation of \$76.6 on a month-level range of
\$121--378.

One design rule, learned from this project's own record and applied to every
day-grain test in Section~\ref{sec:mechanism}: conduct is measured only on
bids formed \emph{outside} direction exposure. On 54\% of essential days the
day-ahead bid was formed while a direction was running or pending, and a bid
formed under direction mechanically shows availability --- treating those days
as choices manufactures state-dependence that evaporates on clean days.
Section~\ref{sec:mechanism} reports the artifact and the correction in full.
